% !TEX root = ../Main.tex
\chapter{逆写系数法与穿根法}

\section{逆写系数法}

遇到首项系数很大、常数项很小的二次式，直接十字相乘不好凑。可以用\textbf{逆写系数}把它换成一个系数更顺手的方程来求根。

\state{逆写系数：倒根等价}{
    方程
    \[
    ax^2+bx+c=0\ (a\ne 0)\quad(\textcircled{1})
    \qquad\text{与}\qquad
    cx^2+bx+a=0\ (c\ne 0)\quad(\textcircled{2})
    \]
    互为\textbf{倒根}。理由：对 $\textcircled{2}$ 两边同除 $x^2$，得 $a\dfrac{1}{x^2}+b\dfrac{1}{x}+c=0$；令 $t=\dfrac{1}{x}$，就变回 $at^2+bt+c=0$，正是 $\textcircled{1}$。所以 $\textcircled{2}$ 的根都是 $\textcircled{1}$ 的根的倒数。
}

\ex{解 $12x^2+x-1>0$。}

\sol{
    把系数逆写：$12x^2+x-1$ 对应方程 $12x^2+x-1=0$，它的倒根方程是 $-x^2+x+12=0$，即
    \[
    x^2-x-12=0\ \Longrightarrow\ (x+3)(x-4)=0,\qquad x=-3,\ 4.
    \]
    取倒数，得 $12x^2+x-1=0$ 的两根 $x=-\dfrac{1}{3},\ \dfrac{1}{4}$，即
    \[
    12x^2+x-1=(3x+1)(4x-1).
    \]
    $a=12>0$，开口向上，$y>0$ 取两根外侧：
    \[
    \text{解集}=\pr{-\infty,-\tfrac{1}{3}}\cup\pr{\tfrac{1}{4},+\infty}.
    \]
}

\rmk{这里逆写出来的 $x^2-x-12$ 正是前面解过的同一个二次式——逆写系数把一个"大系数"的二次不等式拉回到了熟悉的"小系数"二次不等式上。}

\section{穿根法}

二次不等式的图象法可以直接推广到\textbf{高次}：把多项式分解成一串一次因式的乘积，在数轴上一次性读出符号。这就是穿根法（数轴标根法）。

\state{穿根法}{
    记 $y_n=a_0+a_1x+a_2x^2+\cdots+a_nx^n$。若它能写成
    \[
    y_n=a(x-x_1)(x-x_2)\cdots(x-x_m),\qquad x_1<x_2<\cdots<x_m,
    \]
    （其中 $a$ 可以是常数，也可以是常数乘上恒正或恒负的式子，例如 $a=2(x^2+x+1)(x^2+2x+5)$ 恒正），就能用穿根法求 $y_n>0$ 或 $y_n<0$ 的解：

    \begin{enumerate}
        \item 把 $m$ 个根 $x_1,\dots,x_m$ 由小到大标在数轴上；
        \item $a>0$ 从最右端的\textbf{上方}起笔往下穿，$a<0$ 从最右端的\textbf{下方}起笔往上穿；
        \item \textbf{奇重根穿过、偶重根不穿}（碰到偶重根弹回去）；
        \item 曲线在 $x$ 轴上方就是 $y>0$，下方就是 $y<0$。
    \end{enumerate}
}

\begin{center}
\begin{tikzpicture}[>=Stealth, scale=0.82]
    \draw[->] (-0.5,0) -- (6.5,0) node[right] {$x$};
    \foreach \x/\l in {1/x_1, 2.5/x_2, 4/x_3, 5.5/x_m} {
        \filldraw (\x,0) circle (1.2pt);
        \node[below] at (\x, -0.08) {$\l$};
    }
    \draw[thick, blue, smooth] plot coordinates {(0.2,1.4) (1,0) (1.75,-0.7) (2.5,0) (3.25,0.7) (4,0) (4.75,-0.7) (5.5,0) (6.3,1.4)};
    \node[blue] at (3.3, 1.55) {\small $a>0$：从右上起};
\end{tikzpicture}
\qquad
\begin{tikzpicture}[>=Stealth, scale=0.82]
    \draw[->] (-0.5,0) -- (6.5,0) node[right] {$x$};
    \foreach \x/\l in {1/x_1, 2.5/x_2, 4/x_3, 5.5/x_m} {
        \filldraw (\x,0) circle (1.2pt);
        \node[above] at (\x, 0.08) {$\l$};
    }
    \draw[thick, blue, smooth] plot coordinates {(0.2,-1.4) (1,0) (1.75,0.7) (2.5,0) (3.25,-0.7) (4,0) (4.75,0.7) (5.5,0) (6.3,-1.4)};
    \node[blue] at (3.3, -1.55) {\small $a<0$：从右下起};
\end{tikzpicture}
\end{center}

\subsection*{举几个例子}

\ex{$y=(x-2)(x-3)(x-5)$。}

\sol{
    三个单根 $2,3,5$，$a>0$，从右上起、奇穿：
    \begin{center}
    \begin{tikzpicture}[>=Stealth, scale=1]
        \draw[->] (0.5,0) -- (6.5,0) node[right] {$x$};
        \foreach \x in {2,3,5} {\filldraw (\x,0) circle (1.3pt); \node[below] at (\x,-0.08) {$\x$};}
        \draw[thick, blue, smooth] plot coordinates {(0.8,-0.6) (2,0) (2.5,0.5) (3,0) (4,-0.6) (5,0) (6.2,0.9)};
    \end{tikzpicture}
    \end{center}
    $y>0\ \Rightarrow\ 2<x<3$ 或 $x>5$；\qquad $y<0\ \Rightarrow\ x<2$ 或 $3<x<5$。
}

\ex{$y=(x-2)^2(x-5)$。}

\sol{
    根 $x=2$（二重，偶次不穿）、$x=5$（单根），$a>0$：
    \begin{center}
    \begin{tikzpicture}[>=Stealth, scale=1]
        \draw[->] (0.5,0) -- (6.5,0) node[right] {$x$};
        \foreach \x in {2,5} {\filldraw (\x,0) circle (1.3pt); \node[below] at (\x,-0.08) {$\x$};}
        \draw[thick, blue, smooth] plot coordinates {(0.8,-0.8) (1.5,-0.2) (2,0) (2.5,-0.2) (3.5,-0.8) (5,0) (6.2,0.8)};
    \end{tikzpicture}
    \end{center}
    $x=2$ 处弹回去（两侧都在 $x$ 轴下方）。$y>0\ \Rightarrow\ x>5$；\qquad $y<0\ \Rightarrow\ x<2$ 或 $2<x<5$。
}

\ex{$y=(1-x)(x-3)(x-7)$。}

\sol{
    先把每个因式 $x$ 的系数化为正：$1-x=-(x-1)$，于是 $y=-(x-1)(x-3)(x-7)$，\textbf{$a<0$}。三个单根 $1,3,7$，从右下起、奇穿：
    \begin{center}
    \begin{tikzpicture}[>=Stealth, scale=0.95]
        \draw[->] (-0.3,0) -- (8.3,0) node[right] {$x$};
        \foreach \x in {1,3,7} {\filldraw (\x,0) circle (1.3pt); \node[below] at (\x,-0.08) {$\x$};}
        \draw[thick, blue, smooth] plot coordinates {(-0.1,0.8) (1,0) (2,-0.6) (3,0) (5,0.7) (7,0) (8.1,-0.8)};
    \end{tikzpicture}
    \end{center}
    $y>0\ \Rightarrow\ x<1$ 或 $3<x<7$；\qquad $y<0\ \Rightarrow\ 1<x<3$ 或 $x>7$。
}

\rmkb{
    \textbf{方法总结}
    \begin{enumerate}
        \item 判 $a$ 的正负，定穿根的起始方向（$a>0$ 右上、$a<0$ 右下）；
        \item 奇穿偶不穿；
        \item $y$ 正取上、$y$ 负取下。
    \end{enumerate}
}
